\chapter{Análisis del problema}
\label{ch:analisis}

\section{Definición del problema y alcance}
\label{sec:problema-alcance}

\pendiente{definir con precisión el problema que resuelve el sistema y,
especialmente, su alcance: qué entra y qué queda fuera del TFG. Acotar el
alcance es lo que hace evaluable el trabajo.}

\section{Interesados y actores}
\label{sec:actores}

\pendiente{identificar los actores que interactúan con el sistema (corredor,
entrenador o equipo, administrador de datos, etc.), sus objetivos y sus
necesidades principales. Se recomienda una tabla o un diagrama de actores.}

\section{Requisitos funcionales}
\label{sec:requisitos-funcionales}

\pendiente{enumerar los requisitos funcionales, cada uno con identificador
propio y con referencia al objetivo del capítulo \ref{ch:objetivos} que
satisface.}

\section{Requisitos no funcionales}
\label{sec:requisitos-no-funcionales}

\pendiente{enumerar los requisitos no funcionales (rendimiento, usabilidad,
accesibilidad, reproducibilidad, seguridad de credenciales, portabilidad del
almacenamiento, etc.) con criterios medibles siempre que sea posible.}

\section{Requisitos de datos y fuentes}
\label{sec:requisitos-datos}

\pendiente{describir qué datos necesita el sistema, de qué fuentes proceden
(API de Strava, \emph{datasets} simulados, exportaciones propias), con qué
frecuencia y con qué calidad esperada, así como las limitaciones de acceso
(cuotas, autenticación, términos de uso).}

\section{Restricciones}
\label{sec:restricciones}

\pendiente{recoger las restricciones del proyecto: temporales (12 ECTS), de
recursos (equipo local), normativas (Reglamento 24/2024 y guía de desarrollo de
la EPSC) y técnicas (límites de las APIs y de las licencias de datos).}

\section{Criterios de aceptación}
\label{sec:criterios-aceptacion}

\pendiente{definir los criterios que permitirán considerar el sistema aceptado y
verificable, ligados a los requisitos y a los objetivos.}

\section{Catálogo de requisitos}
\label{sec:catalogo-requisitos}

La siguiente tabla es el esqueleto del catálogo de requisitos. Las filas marcadas
como pendientes son ejemplos de marcador de posición y deben sustituirse por
requisitos reales.

\begin{table}[htbp]
  \centering
  \caption{Catálogo de requisitos (esqueleto pendiente de completar).}
  \label{tab:catalogo-requisitos}
  \footnotesize
  \setlength{\tabcolsep}{4pt}
  \begin{tabularx}{\textwidth}{@{}l *{5}{>{\raggedright\arraybackslash}X}@{}}
    \toprule
    \textbf{ID} & \textbf{Descripción} & \textbf{Tipo} & \textbf{Prioridad}
      & \textbf{Origen} & \textbf{Verificación} \\
    \midrule
    \pendiente{RF-001} & \pendiente{descripción del requisito} &
      \pendiente{Funcional} & \pendiente{Alta} & \pendiente{Objetivo 1} &
      \pendiente{prueba que lo verifica} \\
    \addlinespace
    \pendiente{RF-002} & \pendiente{descripción del requisito} &
      \pendiente{Funcional} & \pendiente{Media} & \pendiente{Objetivo 2} &
      \pendiente{prueba que lo verifica} \\
    \addlinespace
    \pendiente{RNF-001} & \pendiente{descripción del requisito} &
      \pendiente{No funcional} & \pendiente{Alta} & \pendiente{Objetivo 4} &
      \pendiente{prueba que lo verifica} \\
    \bottomrule
  \end{tabularx}
\end{table}
